\documentclass[11pt,a4paper]{article}
\usepackage[utf8]{inputenc}\usepackage[T1]{fontenc}\usepackage[spanish]{babel}
\usepackage[margin=2.2cm]{geometry}\usepackage{lmodern}\usepackage{parskip}
\title{Respuesta a la revisión del avance\\\large Detectores YOLO ligeros para plagas y enfermedades del tomate bajo alta radiación solar}
\author{Fray Benjamin Kana Chullo, Hiram Jeremy Ccoto Cruz y Angel Rosendo Condori Coaquira}
\date{17 de noviembre de 2026}
\begin{document}\maketitle
Agradecemos la revisión detallada (reporte G6, 13/11/2026). A continuación se responde cada una de las 17 observaciones y las 4 apreciaciones, con la posición adoptada, lo que se cambió, dónde está en el manuscrito corregido y la evidencia que lo respalda. Posición: \textbf{A} = acepto; \textbf{AP} = acepto en parte; \textbf{P} = pendiente.

El cambio principal es que la radiación solar pasó a ser el aporte central del trabajo (ruta (a) de la observación 1): se midió la exposición, se demostró su efecto con una prueba de estrés controlada y se evaluaron dos estrategias de mitigación en 27 entrenamientos nuevos (v6, con las etiquetas corregidas tras la auditoría y sin la clase hoja sana).

\section*{Resumen}
\begin{itemize}
\item \textbf{Se respondieron las 17 observaciones y las 4 apreciaciones.}
\item \textbf{Atendidas por completo (19):} observaciones 1, 2, 3, 4, 5, 6, 7, 8, 9, 11, 12, 13, 14, 15, 16 y 17; apreciaciones 1, 2, 3 y 4 (en las apreciaciones 3 y 4 se mantiene la decisión, con justificación).
\item \textbf{Obs. 1, 2, 3 y 8:} lo esencial está hecho y medido. Lo que requiere una campaña de campo nueva o más GPU (prueba pareada de luz, experimento controlado de fuga, v4 reentrenada, prueba de la app en campo) queda justificado, con protocolo o script listo y declarado en Limitaciones.
\item \textbf{Obs. 10, en parte:} las reglas, la experiencia de los anotadores, la auditoría y la corrección están hechas y justificadas. El kappa entre anotadores y la validación del ingeniero tienen el protocolo listo (100 y 50 imágenes) y sus resultados se incorporarán al completarse.
\end{itemize}

\section*{Respuesta punto por punto}
\subsection*{Observación 1: Radiación solar no estudiada}
\noindent\textbf{Posición:} \textbf{A} (prueba de campo declarada como trabajo futuro)\\
\textbf{Respuesta:} La radiación pasó a ser el aporte central: título, objetivo, P3 y P4. \textbf{(i) Medida:} al mediodía la saturación se cuadruplica (3.6 \% frente a 0.8–1.0 \%); el Bv del EXIF lo confirma. \textbf{(ii) Efecto causal:} prueba de estrés con la misma foto sobreexpuesta; +1 EV quita 32 \% del mAP y dos tercios del AP de mosca blanca, y la subexposición no afecta. \textbf{(iii) Mitigación:} 27 corridas. El fotométrico eleva la retención a +2 EV (42 \% frente a 29 \%, 9/9 corridas, p = 0.004) a costa de $-$2.4 puntos con luz normal; CLAHE no ayuda. Ninguna estrategia rescata la mosca blanca por encima de +1.5 EV, así que la mitigación debe incluir la captura. \textbf{(iv) Aviso de la app validado:} no salta en ninguna foto real, salta en el 61 \% a +1 EV y el 98 \% desde +1.5 EV, detecta el 79 \% de las fallas inducidas por la luz y, con aviso, el modelo falla el doble (48 \% frente a 24 \%); se mantienen los umbrales y el mensaje ahora dice «Mucha luz: dé sombra». \textbf{Pendiente:} prueba pareada de campo (mismas hojas en 3 franjas, luxómetro, sombra al mediodía, NASA POWER), con el protocolo listo. \textbf{Justificación de lo pendiente:} la prueba pareada de campo necesita una campaña de captura nueva (mismas hojas en tres franjas, con luxómetro) que no cabe en el plazo de este avance. El revisor ofrecía como ruta (a) medir la exposición por imagen, y eso se hizo, junto con una prueba causal más fuerte: la misma foto con distinta exposición. La prueba de campo queda con protocolo escrito y declarada en Limitaciones (iv) y en el trabajo futuro.\\
\textbf{Ubicación en el manuscrito:} Título, resumen, P3–P4, Sec. 3.5, 4.3, 4.4, Tabla 6, Fig. 5, Discusión\\
\textbf{Evidencia:} \texttt{a3\_exposicion.py}, \texttt{a7\_exif\_luz.py}, \texttt{a8\_estres\_exposicion.py}, \texttt{a9\_analisis\_v6.py}, \texttt{a11\_aviso\_luz.py}, \texttt{estres\_v6.csv}, \texttt{aviso\_luz\_resumen.txt}, \texttt{13\_plan\_radiacion.md}, \texttt{15\_resultados\_v6.md}

\subsection*{Observación 2: La caída v1$\rightarrow$v4 no aísla la fuga}
\noindent\textbf{Posición:} \textbf{A} (mínimo aceptable cumplido; experimento controlado declarado)\\
\textbf{Respuesta:} Se reescribió: la caída v1$\rightarrow$v4 ya no se atribuye a la fuga, porque cambiaron a la vez la división, las etiquetas y el formato, y se declaran esos cambios. \textbf{Se cumplió el mínimo aceptable que pide el revisor:} el modelo v1 se evaluó en su propia prueba de 103 imágenes. Con los grupos de captura reconstruidos (203 grupos, 66 repartidos), las 77 imágenes de prueba con foto hermana en entrenamiento o validación dan un mAP@0.5 de \textbf{38.5 \%}, frente a \textbf{12.3 \%} en las 26 sin hermana. Es el mismo modelo con las mismas etiquetas. El conjunto completo da 34.4 \% y la validación 29.0 \%. Se advierte que el subconjunto sin hermana es pequeño y puede diferir en dificultad. \textbf{Justificación de lo pendiente:} el experimento controlado completo (división aleatoria frente a por grupo, con las mismas etiquetas y 3 semillas) está programado y probado en \texttt{entrenar\_control.py} (unas 2.5 h de GPU). No se ejecutó porque la GPU solo estuvo disponible para los 27 entrenamientos de la v6, que tenían prioridad por ser el foco del trabajo. Queda declarado en Limitaciones (v).\\
\textbf{Ubicación en el manuscrito:} Sec. 4.1, 5.1 (v)\\
\textbf{Evidencia:} \texttt{a16\_fuga\_v1.py}, \texttt{v1\_fuga/resumen\_fuga\_v1.txt}, \texttt{entrenar\_control.py}

\subsection*{Observación 3: Comparación v4–v5 confundida}
\noindent\textbf{Posición:} \textbf{A}\\
\textbf{Respuesta:} (i) Confirmado: los valores v4 del «test fijo» eran los de 102 imágenes con etiquetas v4. Se reevaluó todo con las mismas 97 imágenes y las etiquetas v5. (ii) Se declara la versión de etiquetas. (iii) Se reconocen los cambios de hiperparámetros. (iv) La conclusión cambió: la v5 no mejora en mayo y sí en julio (generalización temporal). La v6 confirma el patrón (prueba fija 20.2 \%, nuevas 37.6 \%). El control «solo mayo frente a mixto del mismo tamaño», que separa sesiones de cantidad, está en \texttt{entrenar\_control.py} sin ejecutar. \textbf{Justificación de lo pendiente:} el reentrenamiento de la v4 con la configuración de la v5 no se ejecutó por tiempo de GPU. Por eso la conclusión se formuló con cautela («debe confirmarse») y se declara en Discusión y en Limitaciones (v). El modelo \texttt{s} no se añadió porque el alcance se acota a modelos nano desplegables en el teléfono (apreciación 4).\\
\textbf{Ubicación en el manuscrito:} Tabla 5, Fig. 2, Sec. 4.2, 4.4, Discusión\\
\textbf{Evidencia:} \texttt{metricas/map\_ultralytics.csv}, \texttt{v4\_vs\_v5\_pareado.csv}, \texttt{comparacion\_v6.csv}

\subsection*{Observación 4: Texto en español para Sensors}
\noindent\textbf{Posición:} \textbf{A}\\
\textbf{Respuesta:} Se preparó la versión completa en inglés (\texttt{main\_en.tex}, \texttt{Article\_revised\_MDPI\_KanaChullo\_EN.pdf}, 14 páginas). Tiene el mismo contenido, tablas y cifras que la versión en español, con el resumen en 170 palabras y terminología consistente (*late blight*, *powdery mildew*, *whitefly*, *tomato leafminer*, *false-healthy rate*, *capture group*). Las figuras se regeneraron con rótulos en inglés a partir de los mismos datos. Las capturas de la app quedan en español y así se indica en el pie de figura. Se recomienda una revisión de estilo por un hablante nativo antes del envío a la revista. Cada cambio futuro de fondo se aplicará a ambas versiones.\\
\textbf{Ubicación en el manuscrito:} Manuscrito completo (versión EN)\\
\textbf{Evidencia:} \texttt{main\_en.tex}, \texttt{a15\_figuras\_en.py}

\subsection*{Observación 5: Grupo por tiempo, no por planta}
\noindent\textbf{Posición:} \textbf{A}\\
\textbf{Respuesta:} Se usa «grupo de captura» en lugar de «planta». Auditoría independiente del tiempo con pHash sobre los 856 785 pares entre subconjuntos: 0 pares casi idénticos (mínimo 12; referencia intra-ráfaga 30). Se declara que no descarta la misma planta en dos ráfagas. «Imágenes nuevas» (julio) funciona como prueba temporal parcial. La prueba de campo nueva se dividirá por planta marcada.\\
\textbf{Ubicación en el manuscrito:} Sec. 3.3, 4.1, 5.1\\
\textbf{Evidencia:} \texttt{duplicados/resumen.json}, \texttt{a4\_duplicados.py}, \texttt{13\_plan\_radiacion.md}

\subsection*{Observación 6: Solo DE entre semillas}
\noindent\textbf{Posición:} \textbf{A}\\
\textbf{Respuesta:} IC 95 \% por bootstrap agrupado (1000 réplicas): unos $\pm$6–8 puntos de mAP@0.5 en la prueba fija. Diferencias pareadas v4–v5 con IC; McNemar e IC de Wilson; n de cajas por clase. En la v6, las diferencias entre experimentos se prueban por pares de arquitectura y semilla (Wilcoxon) y se informa cuántos pares mejoran. Se declara que con 3 semillas la DE es del orden de varias diferencias.\\
\textbf{Ubicación en el manuscrito:} Sec. 3.5, 4.2, 4.4, 4.5, Tablas 5–7, 5.1 (vi)\\
\textbf{Evidencia:} \texttt{map\_con\_ic.csv}, \texttt{mcnemar\_v4\_v5.csv}, \texttt{a9\_analisis\_v6.py}, \texttt{robustez\_v6.csv}

\subsection*{Observación 7: Acierto de clase principal mal definido}
\noindent\textbf{Posición:} \textbf{A}\\
\textbf{Respuesta:} Métricas por imagen: sensibilidad por plaga y tasa de falso sano (22 \% v4; 28 \% v5 instalado; 44 \% v6 base, 33–60 \% entre corridas). Líneas base «siempre mildiu» (60.8 \% en la prueba fija, igual que los modelos) y «siempre sana». El falso sano se mide también en la prueba de estrés: sube de 44 \% a 65 \% con +2 EV. La discusión propone que la app muestre la confianza y pida revisar cuando no detecta nada. Falta calibración y regla de abstención.\\
\textbf{Ubicación en el manuscrito:} Sec. 3.5, 4.3, 4.5, Tablas 6–7, Discusión\\
\textbf{Evidencia:} \texttt{metricas\_por\_imagen.csv}, \texttt{lineas\_base.csv}, \texttt{estres\_v6.csv}

\subsection*{Observación 8: Teléfono con modelos v1/v4; discordancia}
\noindent\textbf{Posición:} \textbf{A} (campo P)\\
\textbf{Respuesta:} v5 convertido (6.1 MB; coincide con ONNX en 96/96, $\Delta$ $\leq$ 0.014) e instalado. La discordancia se explica por la interpolación del redimensionado (OpenCV frente a Dart). Se reporta el delegado (CPU, 4 hilos). Las capturas se declaran tomadas sobre un monitor. \textbf{v6 desplegado:} backend y app adaptados a 4 clases con el YOLOv8n v6 (servidor \texttt{best.pt}; TFLite float16 de 6.1 MB, coincide con ONNX en 96/96, $\Delta$ = 0.004; pruebas del backend 14/14 y app sin errores de análisis). \textbf{Medido en el Galaxy A56 con la v6} (192 imágenes, el mismo código de la app): 574 $\pm$ 186 ms por foto (mediana 596 ms; primera corrida 756 $\pm$ 497 ms; detecciones idénticas en ambas). El tiempo lo domina decodificar y reducir la foto: 271 ms con fotos < 1 MP frente a 678 ms con fotos de 50 MP. Misma clase principal que el servidor en el 82.8 \% (IC 95 \%: 76.8–87.5), igual que la v4. Reproduce el 77 \% de las cajas del servidor. \textbf{Pendiente:} igualar la interpolación y probar en campo.\\
\textbf{Ubicación en el manuscrito:} Sec. 3.6, 4.5, Tabla 3, Fig. 6, 5.1 (vii)\\
\textbf{Evidencia:} \texttt{comparar\_onnx\_tflite.py}, \texttt{v6/celular/comparacion\_servidor\_movil\_v6.json}

\subsection*{Observación 9: Sin trabajos relacionados}
\noindent\textbf{Posición:} \textbf{A}\\
\textbf{Respuesta:} Nueva sección con tabla. Solo se citan trabajos leídos en la fuente: Shehu 2025, Debnath 2023, Zhao 2025, Gerz y Jelali 2025. Se precisa el vacío: ninguno mide cómo la sobreexposición degrada al detector. Kamei 2023, Gookyi 2024 y HFA-Net no se pudieron leer (acceso bloqueado).\\
\textbf{Ubicación en el manuscrito:} Sec. 2, Tabla 1\\
\textbf{Evidencia:} —

\subsection*{Observación 10: Etiquetas sin validez experta}
\noindent\textbf{Posición:} \textbf{A} en reglas y justificación; \textbf{AP} en el kappa (resultados en curso)\\
\textbf{Respuesta:} \textbf{(i) Experiencia declarada:} el anotador principal tiene 7 años en el cultivo de tomate en El Pedregal (Majes, Arequipa), en un fundo con asistencia técnica de ingenieros de campo. Se declara como experiencia de campo, no como sustituto de una validación independiente. \textbf{(ii) Reglas y alcance:} reglas escritas completas (incluido 16 px); especie de oídio no confirmada; polilla como insecto + mina; identificación por síntomas visibles. \textbf{(ii-b) Reglas completas ahora en el texto:} la regla de 16 px se aplicó a las lesiones. La mosca blanca se etiquetó aunque quedara por debajo de ese tamaño (el 87 \% de sus cajas) porque es la plaga de interés y omitirla dejaría la clase casi vacía. Eso explica su bajo AP, y para la versión siguiente se propone una caja por grupo de insectos. La polilla une insecto y mina porque la mina es el daño visible de la misma plaga, la acción de manejo es la misma y, con 473 cajas, separarla dejaría dos clases demasiado pequeñas. \textbf{(iii) Auditoría y corrección:} 9 modelos como revisores encontraron 4 problemas. En la v6 se revisaron a mano 319 imágenes (mildiu 2698 $\rightarrow$ 2933 cajas) y se eliminó la hoja sana. Solo 2 cajas tenían la clase equivocada: el problema es de criterio, no de identificación. \textbf{(iv) Acuerdo entre anotadores (protocolo ciego, en curso):} muestra estratificada de 100 imágenes (3 sesiones, un grupo por imagen, 258 cajas), etiquetada por separado por los dos anotadores. Se reporta el kappa de Cohen por clase y el F1 por caja (IoU $\geq$ 0.5 y $\geq$ 0.3). El script está validado. \textbf{(v) Validación externa (en curso):} hoja de 50 fotos para un ingeniero agrónomo; kappa con las etiquetas v6. \textbf{Justificación de lo pendiente:} el protocolo sigue al pie de la letra la recomendación del revisor (muestra estratificada de 100 imágenes, segundo anotador, kappa y acuerdo de cajas con IoU $\geq$ 0.5). Requiere unas 3 h de etiquetado ciego de cada anotador y la disponibilidad de un ingeniero externo, que no caben en el plazo de esta entrega. Los scripts de cálculo están validados y las cifras se incorporarán en cuanto lleguen las etiquetas.\\
\textbf{Ubicación en el manuscrito:} Sec. 3.2, 3.5 (acuerdo), 4.6, Tabla 2, Fig. 7, 5.1 (i)\\
\textbf{Evidencia:} \texttt{errores/candidatos.csv}, \texttt{12\_correccion\_etiquetas.md}, \texttt{a12\_muestra\_acuerdo.py}, \texttt{a13\_acuerdo\_anotadores.py}, \texttt{a14\_hoja\_ingeniero.py}, \texttt{acuerdo\_anotadores/}

\subsection*{Observación 11: Datos y código no publicados}
\noindent\textbf{Posición:} \textbf{A}\\
\textbf{Respuesta:} Dataset público en Roboflow Universe (CC BY 4.0): versión 1 = v4, versión 2 = v5, versión 3 = v6. Código, listas de prueba, métricas y modelos en GitHub; tabla de entorno con versiones. Opcional: depósito en Zenodo con DOI.\\
\textbf{Ubicación en el manuscrito:} Sec. 3.7, Tablas 3–4, Disponibilidad\\
\textbf{Evidencia:} Links verificados

\subsection*{Observación 12: Formato de informe de avance}
\noindent\textbf{Posición:} \textbf{A}\\
\textbf{Respuesta:} Objetivo y preguntas P1–P5; resultados organizados por pregunta; sin «PPI», fases ni versión del servidor; párrafo de limitaciones.\\
\textbf{Ubicación en el manuscrito:} Introducción, Sec. 4, 5.1\\
\textbf{Evidencia:} —

\subsection*{Observación 13: Figuras}
\noindent\textbf{Posición:} \textbf{A}\\
\textbf{Respuesta:} Se eliminaron el diagrama de fases, el mosaico de Ultralytics y la figura ilegible de curvas. Figuras nuevas: comparación con las mismas etiquetas, por clase, exposición, prueba de estrés (con banda de $\pm$1 DE) y auditoría.\\
\textbf{Ubicación en el manuscrito:} Figs. 1–7\\
\textbf{Evidencia:} —

\subsection*{Observación 14: Cifras sin trazabilidad}
\noindent\textbf{Posición:} \textbf{A}\\
\textbf{Respuesta:} «Hasta 6 puntos» se reemplazó por IC y rangos. «96 de 96» se explica como la comparación ONNX–TFLite con la misma entrada. Se retiró «más GPU no ayudan». Cada cifra de la v6 sale de \texttt{comparacion\_v6.csv} / \texttt{estres\_v6.csv}.\\
\textbf{Ubicación en el manuscrito:} Sec. 4.2, 4.4, 4.5\\
\textbf{Evidencia:} CSV citados

\subsection*{Observación 15: Referencias incompletas}
\noindent\textbf{Posición:} \textbf{A}\\
\textbf{Respuesta:} Se completó [JutePest] con volumen, páginas y DOI (IEEE Access 12: 72938–72956, doi 10.1109/ACCESS.2024.3403491, verificado en Crossref). Se retiraron las referencias débiles o sin verificar ([1] guía institucional, [5], [8], [11], [12] y otras: SENASA, Shafi, Ksibi, Das, Mengesha, Yadav, Feng). Todas las referencias de revista quedan con DOI en formato MDPI. El software (Ultralytics) y el dataset se citan con su URL estable.\\
\textbf{Ubicación en el manuscrito:} Referencias\\
\textbf{Evidencia:} —

\subsection*{Observación 16: Terminología}
\noindent\textbf{Posición:} \textbf{A}\\
\textbf{Respuesta:} Nombres comunes y científicos en el texto; sin identificadores de código.\\
\textbf{Ubicación en el manuscrito:} Todo el texto\\
\textbf{Evidencia:} —

\subsection*{Observación 17: Título, resumen y palabras clave}
\noindent\textbf{Posición:} \textbf{A}\\
\textbf{Respuesta:} Título nuevo centrado en la radiación. Resumen de menos de 200 palabras con objetivo, método y resultados cuantitativos. Palabras clave que no repiten el título (radiación solar, sobreexposición, aumento fotométrico, etc.).\\
\textbf{Ubicación en el manuscrito:} Portada\\
\textbf{Evidencia:} —

\subsection*{Apreciación 1: Hoja sana como clase}
\noindent\textbf{Posición:} \textbf{A} (cambio)\\
\textbf{Respuesta:} Se aceptó la sugerencia. En la v5 se reporta el mAP con y sin ella (baja 6–9 puntos). En la v6 se \textbf{eliminó}: estaba etiquetada en parte (212 casos sin caja) y no es una plaga. Las imágenes sin plaga quedan como fondo, y el falso sano mide el error de declarar sana una planta afectada.\\
\textbf{Ubicación en el manuscrito:} Sec. 3.2, Tablas 2 y 5–6\\
\textbf{Evidencia:} \texttt{map\_4\_vs\_5.csv}, \texttt{comparacion\_v6.csv}

\subsection*{Apreciación 2: Conservar la etapa v1}
\noindent\textbf{Posición:} \textbf{A}\\
\textbf{Respuesta:} Se conservó la etapa v1 reducida a un párrafo en P1, como caso medido de fuga. Ahora además aporta la medición directa: el mismo modelo v1 rinde 38.5 \% en las imágenes de prueba con foto hermana frente a 12.3 \% en las que no la tienen.\\
\textbf{Ubicación en el manuscrito:} Sec. 4.1\\
\textbf{Evidencia:} —

\subsection*{Apreciación 3: Elegir YOLOv8n aunque YOLOv10n gane en prueba}
\noindent\textbf{Posición:} \textbf{Mantengo}\\
\textbf{Respuesta:} Seleccionar con prueba sesga la estimación (Kapoor y Narayanan). En la v6, YOLOv8n gana en validación y también en prueba.\\
\textbf{Ubicación en el manuscrito:} Discusión\\
\textbf{Evidencia:} \texttt{resumen\_v6.txt}

\subsection*{Apreciación 4: Solo modelos nano}
\noindent\textbf{Posición:} \textbf{Mantengo}\\
\textbf{Respuesta:} Restricción de despliegue en el teléfono; la conclusión se acota a modelos ligeros. El modelo \texttt{s} queda como experimento opcional.\\
\textbf{Ubicación en el manuscrito:} Discusión\\
\textbf{Evidencia:} —

\section*{Pendientes antes del envío a la revista}
\begin{enumerate}
\item Prueba pareada de campo de la radiación (obs. 1), según el protocolo de \texttt{13\_plan\_radiacion.md}.
\item Ejecutar \texttt{entrenar\_control.py} en la RTX, unas 2.5 h (obs. 2 y 3), para complementar la medición de la v1.
\item Etiquetado ciego de las 100 imágenes por Fray y Hiram, y hoja de 50 fotos para un ingeniero agrónomo (obs. 10).
\item Igualar la interpolación servidor–teléfono y probar la app en campo (obs. 8).
\item Opcional: leer y citar Kamei 2023, Gookyi 2024 y HFA-Net si se consigue acceso (obs. 9).
\item Revisión de estilo del inglés por un hablante nativo y aplicar a la versión EN cualquier cambio futuro (obs. 4).
\end{enumerate}
\end{document}
